\documentclass[12pt,reqno]{amsart}

\usepackage{amsmath,amssymb,amsthm}
\usepackage{booktabs}
\usepackage[colorlinks=true,linkcolor=blue,citecolor=blue,urlcolor=blue]{hyperref}
\usepackage[margin=1.15in]{geometry}
\usepackage{enumitem}
\usepackage{microtype}

\newtheorem{theorem}{Theorem}[section]
\newtheorem{proposition}[theorem]{Proposition}
\newtheorem{corollary}[theorem]{Corollary}
\theoremstyle{definition}
\newtheorem{definition}[theorem]{Definition}
\theoremstyle{remark}
\newtheorem{remark}[theorem]{Remark}
\newtheorem{question}[theorem]{Question}

\newcommand{\ZZ}{\mathbb{Z}}
\newcommand{\QQ}{\mathbb{Q}}
\newcommand{\Gstar}{G^{\!*}}
\newcommand{\varpiup}{\varpi}
\newcommand{\agm}{\operatorname{agm}}
\newcommand{\abs}[1]{\lvert #1 \rvert}

\title{The Missing Ratio:\\[2pt]
  $\Gamma(\tfrac14)/\Gamma(\tfrac34)$ as a Fundamental Constant}

\author{}
\address{}

\subjclass[2020]{33B15, 11M06, 01A55, 33E05}
\keywords{Gamma function, reflection formula, lemniscate constant,
  Euler, fundamental constants, Wallis product}

\date{\today}

\begin{document}

\begin{abstract}
The Euler reflection formula at $z=\tfrac14$ gives
$\Gamma(\tfrac14)\,\Gamma(\tfrac34) = \pi\sqrt{2}$.
This \emph{product} is classical and well studied.  We observe that the
\emph{ratio} $\Gamma(\tfrac14)/\Gamma(\tfrac34) \approx 2.9587$---which
we denote $\Gstar$ and call the \emph{lemniscatic bridge constant}---has
received almost no independent attention, despite admitting a Wallis-type
infinite product, appearing as the scaled period of the CM elliptic curve
$y^{2}=x^{3}-x$, and factoring the lemniscate constant as
$\varpiup = {\Gstar}\sqrt{\pi}/2$.

We argue that $\Gstar$ deserves recognition as a fundamental constant on
the same footing as $\pi$ and $\varpiup$, and that its historical
absence from the literature is an accident of convention, not of
substance.
\end{abstract}

\maketitle
\tableofcontents


% ================================================================
\section{The reflection formula produces two objects}
% ================================================================

The Euler reflection formula states that for all $z\notin\ZZ$,
\begin{equation}\label{eq:reflection}
  \Gamma(z)\,\Gamma(1-z) \;=\; \frac{\pi}{\sin(\pi z)}\,.
\end{equation}
At $z = \tfrac14$:
\begin{equation}\label{eq:product}
  \Gamma(\tfrac14)\,\Gamma(\tfrac34)
  \;=\; \frac{\pi}{\sin(\pi/4)}
  \;=\; \pi\sqrt{2}\,.
\end{equation}
This is the \emph{product}.  It yields a circular quantity ($\pi$,
modified by $\sqrt{2}$) and has been known since Euler.  It appears in
hundreds of papers and textbooks.

The same two Gamma values also determine a \emph{ratio}:
\begin{equation}\label{eq:ratio}
  \frac{\Gamma(\tfrac14)}{\Gamma(\tfrac34)}
  \;=\; \frac{\Gamma(\tfrac14)^{2}}{\pi\sqrt{2}}
  \;=\; 2.9586751191886388\ldots
\end{equation}
This number, which we denote $\Gstar$, has---to our knowledge---never
been given a standard name or symbol, despite being as elementary as
the product.

\begin{remark}
Every evaluation of the reflection formula at a rational argument
$z = p/q$ produces a product and a ratio.  At $z = \tfrac12$, the
product $\Gamma(\tfrac12)^{2} = \pi$ yields the circle constant itself,
and the ratio $\Gamma(\tfrac12)/\Gamma(\tfrac12) = 1$ is trivial.
At $z = \tfrac14$, the product yields $\pi\sqrt{2}$ and the ratio
yields $\Gstar$.  This is the \emph{first} non-trivial ratio produced
by the reflection formula at a point of the form $z = 1/(2n)$.
\end{remark}


% ================================================================
\section{Why the ratio was overlooked}
% ================================================================

We identify three historical reasons why $\Gstar$ was not recognised as
a fundamental constant.

\subsection{The tradition of periods}

Since Euler, the preferred objects in elliptic function theory have been
\emph{periods}---the arc lengths, areas, and integrals that arise as
lattice generators.  The lemniscate constant
\begin{equation}\label{eq:varpi}
  \varpiup \;=\; \frac{\Gamma(\tfrac14)^{2}}{2\sqrt{2}\,\Gamma(\tfrac12)}
  \;=\; 2.6220575542921120\ldots
\end{equation}
is the half-perimeter of the lemniscate $r^{2}=\cos 2\theta$ and a
period of the elliptic curve $y^{2}=x^{3}-x$.  As a period, it has
geometric meaning: it is a length one can draw.  The constant $\Gstar$
is not directly a period; it is a \emph{ratio} of two Gamma values.
Periods were named and studied; ratios were computed and discarded.

\subsection{The dominance of $\pi$}

The reflection formula~\eqref{eq:product} was primarily used to compute
$\pi$ and its relatives, not to study the Gamma function's internal
structure.  The product $\Gamma(\tfrac14)\Gamma(\tfrac34) = \pi\sqrt{2}$
was valued because it gave $\pi$; the ratio was an intermediate quantity
with no known application.

\subsection{The $\pi$-contaminated notation}

The standard representation
$\varpiup = \Gamma(\tfrac14)^{2}/(2\sqrt{2\pi})$ merges $\Gamma(\tfrac14)$
and $\pi$ into a single expression, hiding the factorisation
$\varpiup = \Gstar\sqrt{\pi}/2$.  In this notation, $\Gstar$ is invisible:
it appears only as a ratio of two subexpressions of $\varpiup$ that are
never separated.  The notation discouraged the decomposition.


% ================================================================
\section{The constant and its properties}
% ================================================================

\begin{definition}\label{def:Gstar}
The \emph{lemniscatic bridge constant} is
\begin{equation}\label{eq:Gstar}
  {\Gstar} \;=\; \frac{\Gamma(\tfrac14)}{\Gamma(\tfrac34)}
  \;=\; 2.9586751191886388\ldots
\end{equation}
\end{definition}

We collect its principal representations, all classical:
\begin{align}
  {\Gstar} &= \frac{\Gamma(\tfrac14)^{2}}{\sqrt{2}\,\Gamma(\tfrac12)^{2}}
    &&\text{($\pi$-free form)}\label{eq:pifree}\\[4pt]
  {\Gstar} &= \frac{2\varpiup}{\sqrt{\pi}}
    &&\text{(lemniscate form)}\label{eq:lem}\\[4pt]
  {\Gstar} &= \frac{2\sqrt{2}}{\sqrt{\pi}}\,
    K\!\bigl(\tfrac{1}{\sqrt{2}}\bigr)
    &&\text{(elliptic integral)}\label{eq:elliptic}\\[4pt]
  {\Gstar} &= \frac{2\sqrt{\pi}}{\agm(1,\sqrt{2})}
    &&\text{(AGM)}\label{eq:agm}\\[4pt]
  {\Gstar} &= \sqrt{2\pi}\;\theta_{3}(e^{-\pi})^{2}
    &&\text{(theta function at self-dual nome)}\label{eq:theta}\\[4pt]
  {\Gstar} &= \frac{8\,L(E,1)}{\sqrt{\pi}}
    &&\text{($L$-function of $y^{2}=x^{3}-x$)}\label{eq:Lfunction}
\end{align}

Each representation is known; none was previously attributed to a named
constant.


% ================================================================
\section{What the ratio knows that the product does not}
% ================================================================

The product $\Gamma(\tfrac14)\Gamma(\tfrac34)$ is \emph{symmetric}: it
treats $z=\tfrac14$ and $z=\tfrac34$ on equal footing.  Swapping them
changes nothing.  The product yields $\pi\sqrt{2}$, a quantity that
depends only on the \emph{sum} $z + (1-z) = 1$, not on the specific
partition into $\tfrac14$ and $\tfrac34$.

The ratio $\Gamma(\tfrac14)/\Gamma(\tfrac34)$ is \emph{asymmetric}: it
distinguishes $z=\tfrac14$ from $z=\tfrac34$.  It measures \emph{how much
larger} $\Gamma(\tfrac14)$ is than $\Gamma(\tfrac34)$.  This asymmetry
carries information:

\begin{theorem}[Wallis product for the ratio]\label{thm:wallis}
\begin{equation}\label{eq:wallis-ratio}
  {\Gstar} \;=\; \lim_{N\to\infty}\;(N+1)^{-1/2}
  \prod_{k=0}^{N}\frac{4k+3}{4k+1}\,.
\end{equation}
\end{theorem}

\begin{proof}
Write $Q_{N} = \prod_{k=0}^{N}(4k+3)/(4k+1)
= (\tfrac34)_{N+1}/(\tfrac14)_{N+1}
= \Gamma(\tfrac34+N+1)\Gamma(\tfrac14)/
[\Gamma(\tfrac14+N+1)\Gamma(\tfrac34)]$.
By the asymptotic ratio $\Gamma(x+a)/\Gamma(x+b)\sim x^{a-b}$,
$Q_{N}\sim [\Gamma(\tfrac14)/\Gamma(\tfrac34)]\,(N+1)^{1/2}$.
\end{proof}

The numerators $3, 7, 11, 15, 19, 23, \ldots$ are integers
$\equiv 3\pmod{4}$.  The denominators $1, 5, 9, 13, 17, 21, \ldots$
are integers $\equiv 1\pmod{4}$.  In the arithmetic of the Gaussian
integers $\ZZ[i]$, the former are \emph{inert} (they remain prime) and
the latter include the \emph{split} primes (they factor as
$\pi\bar{\pi}$).  The ratio measures the cumulative advantage of
inertness over splitting---the Fermat two-square classification made
multiplicative.

The product knows about $\pi$.  The ratio knows about the lattice.


% ================================================================
\section{The triad and the composite nature of $\varpiup$}
% ================================================================

\begin{theorem}[Factorisation of $\varpiup$]\label{thm:factor}
\begin{equation}\label{eq:factor}
  \varpiup \;=\; \frac{{\Gstar}\cdot\sqrt{\pi}}{2}\,.
\end{equation}
\end{theorem}

\begin{proof}
$\Gstar\sqrt{\pi}/2
= [\Gamma(\tfrac14)/\Gamma(\tfrac34)]\cdot\Gamma(\tfrac12)/2$.
By reflection, $\Gamma(\tfrac34)=\pi\sqrt{2}/\Gamma(\tfrac14)$, so
$\Gstar\sqrt{\pi}/2
= \Gamma(\tfrac14)^{2}\Gamma(\tfrac12)/(2\pi\sqrt{2})
= \Gamma(\tfrac14)^{2}/(2\sqrt{2}\Gamma(\tfrac12)) = \varpiup$.
\end{proof}

The lemniscate constant, treated since Euler as a single irreducible
period, factors into two independent pieces:
\begin{itemize}
  \item $\sqrt{\pi}$, which arises from the Wallis product over the
    even/odd partition (the mod-$2$ ``race'');
  \item $\Gstar$, which arises from the Wallis product over the
    $1\bmod 4$ / $3\bmod 4$ partition (the mod-$4$ ``race'').
\end{itemize}
Neither piece alone yields $\varpiup$; both are required.

\begin{proposition}[No single-product representation of $\varpiup$]
\label{prop:no-third}
There is no ratio product $\prod(an+b)/(an+c)$ over a single arithmetic
progression that converges (after polynomial prefactor) to $\varpiup$.
\end{proposition}

\begin{proof}[Argument]
Such a product converges to a ratio $\Gamma(c/a)/\Gamma(b/a)$ (times
algebraic factors), which involves $\Gamma(\tfrac14)$ at most once.
The lemniscate constant requires $\Gamma(\tfrac14)^{2}$ in the
numerator, which cannot be produced by a single ratio.
\end{proof}


% ================================================================
\section{The triad identity}
% ================================================================

\begin{theorem}\label{thm:triad}
$\pi = 4\varpiup^{2}/{\Gstar}^{2}$.
\end{theorem}

\begin{proof}
$4\varpiup^{2}/{\Gstar}^{2}
= 4\cdot({\Gstar}\sqrt{\pi}/2)^{2}/{\Gstar}^{2}
= 4\cdot{\Gstar}^{2}\pi/(4{\Gstar}^{2}) = \pi$.
\end{proof}

This expresses the circle constant as a ratio of two lemniscatic
objects.  In the hierarchy
$\Gamma(\tfrac14) \to \Gstar \to \varpiup \to \pi$,
the ratio precedes the period, and the period precedes the circle.


% ================================================================
\section{Why this matters}
% ================================================================

The constant $\Gstar = \Gamma(\tfrac14)/\Gamma(\tfrac34)$ appears
identically in at least seven branches of mathematics (see~\cite{Paper2}
for a catalogue of fifty-two representations):
\begin{center}
\begin{tabular}{ll}
\toprule
\textbf{Branch} & \textbf{Form} \\
\midrule
Complex analysis & Gamma ratio $\Gamma(\tfrac14)/\Gamma(\tfrac34)$ \\
Classical geometry & Scaled lemniscate arc length $2\varpiup/\sqrt{\pi}$ \\
Elliptic integrals & $2\sqrt{2}\,K(1/\!\sqrt{2})/\sqrt{\pi}$ \\
Computational number theory & $2\sqrt{\pi}/\agm(1,\sqrt{2})$ \\
Modular forms & $\sqrt{2\pi}\,\theta_{3}(e^{-\pi})^{2}$ \\
Lattice combinatorics & $\sqrt{2\pi\,W_{3}}$ (Watson integral) \\
Arithmetic geometry & $8\,L(E,1)/\sqrt{\pi}$ (BSD formula) \\
\bottomrule
\end{tabular}
\end{center}

No other unnamed constant connects this many domains.  The constant
$\pi$ is the product side of the reflection formula; $\Gstar$ is the
ratio side.  Both are outputs of the same algebraic identity at the
same evaluation point.  One was elevated to the most famous number
in mathematics.  The other was left on the floor.

The oversight is understandable.  The product gives $\pi$, which has
obvious geometric meaning (circumference, area, periodicity).  The
ratio gives $\Gstar$, which has no simple geometric interpretation as
a length or an area.  But $\Gstar$ has \emph{arithmetic} meaning: it
is the limit of the Wallis product over the Fermat residue classes,
the scaled period of the simplest CM curve, the bridge between the
lattice and the circle.

The question is not whether $\Gstar$ is new---it is not.
The question is whether it is fundamental.  We believe the evidence
presented here, and in the companion papers~\cite{Paper2,Paper3},
answers in the affirmative.


% ================================================================
\section{A note on naming}
% ================================================================

We have used the notation $\Gstar$ throughout, following its role as the
``bridge constant'' in Foundational Ternary Dynamics.  We note that the
constant admits a simpler description than any notation suggests:

\medskip
\begin{center}
\emph{$\Gstar$ is the ratio of the Gamma function at the
quarter-point to its reflection.}
\end{center}
\medskip

\noindent
A dedicated symbol is warranted by the same criterion that warranted
$\pi$ and $\varpiup$: the constant appears independently in enough
contexts that writing it out in full each time obscures the structure
it reveals.


% ================================================================
% REFERENCES
% ================================================================

\begin{thebibliography}{10}

\bibitem{BorweinBorwein}
J.~M.~Borwein and P.~B.~Borwein,
\emph{Pi and the AGM}, Wiley, 1987.

\bibitem{Finch}
S.~R.~Finch,
\emph{Mathematical Constants}, Cambridge Univ.\ Press, 2003.

\bibitem{Paper2}
\emph{Fifty-two faces of the lemniscatic bridge constant},
companion paper.

\bibitem{Paper3}
\emph{The two races: Wallis products for $\sqrt{\pi}$, $\Gstar$, and
the composite nature of $\varpiup$},
companion paper.

\bibitem{Wallis}
J.~Wallis,
\emph{Arithmetica Infinitorum}, Oxford, 1656.

\end{thebibliography}

\end{document}
